% Scheletro di lavoro, non testo definitivo.
% Fonti e passi sono proposte di lettura: verificarli sul testo greco prima della stesura.
% Seguira il metodo: documento originale, traduzione, esame dell'argomentazione, ricostruzione critica.
\chapter{La conoscenza e la dialettica}
\label{cap:platone:conoscenza}

% Problema: È possibile cercare ciò che non si conosce?
% Fonti da verificare: Menone 80d–86c, 97a–100b.
\section{Il problema dell'apprendimento}

\subsection{Il paradosso della ricerca}

\subsection{La reminiscenza e l'esempio geometrico}

% Problema: Perché sensazione e opinione vera non bastano a definire la conoscenza?
% Fonti da verificare: Teeteto 151d–210d.
\section{Che cosa significa sapere?}

\subsection{Conoscenza e percezione}

\subsection{Opinione vera, logos e aporia finale}

% Problema: Come distinguere i gradi del sapere e gli oggetti corrispondenti?
% Fonti da verificare: Repubblica V 476a–480a; VI 509d–511e.
\section{Le forme della conoscenza}

\subsection{Doxa ed episteme}

\subsection{La linea divisa e il ruolo della matematica}

% Problema: Che cosa distingue l'indagine dialettica dalla semplice disputa?
% Fonti da verificare: Repubblica VI–VII; Fedro 265d–266c; Sofista 253b–e.
\section{Il procedimento dialettico}

\subsection{Dall'ipotesi al principio}

\subsection{Raccolta, divisione e connessione dei generi}
